\section{Discussion}

Ce projet s'est avéré beaucoup plus difficile que prévu, présentant des obstacles tant sur le plan théorique que méthodologique et technique.

La lecture et la compréhension des articles ont été difficiles. La théorie des processus stochastiques nécessaire pour comprendre les GPIS dépasse certains aspects de mon bagage mathématique, m'imposant de revenir sur des notions fondamentales avant d'aborder le matériel avancé. En tant que première expérience de recherche formelle, j'ai d'abord eu des difficultés à analyser efficacement les articles; j'ai adopté une stratégie en deux passes : une lecture initiale pour saisir la structure globale, puis une seconde passe détaillée où je condense chaque paragraphe en une phrase en utilisant Obsidian pour organiser mes notes.

Le développement de la preuve de concept pour Mitsuba 3 a présenté des défis techniques importants. L'environnement de compilation a exigé une correspondance précise des versions (GCC, CMake, Ninja, Embree). De plus, la documentation limitée de l'API C++ de Mitsuba 3 a nécessité une rétro-ingénierie par inspection du code, et la maîtrise de la syntaxe DrJit a ajouté une couche de complexité supplémentaire.

Malgré ces difficultés, ce projet a été intellectuellement stimulant et professionnellement enrichissant. J'ai acquis une exposition précieuse aux développement en recherche en infographie, notamment l'intégration de techniques d'apprentissage automatique pour améliorer le photoréalisme et les applications pratiques. Le projet m'a aussi obligé à revisiter des concepts fondamentaux du tracé de rayon et des fonctions de diffusion, renforçant ma compréhension des sujets avancés au coeur de ce travail. Surtout, cette expérience a développé ma capacité à aborder des concepts théoriques complexes et à les appliquer à des bases de code de production existantes, compétence cruciale pour la recherche et le développement en infographie.

Bien que le projet ait été plus exigeant que prévu, j'anticipe des résultats significatifs de l'intégration des GPIS dans Mitsuba 3. Une part importante de théorie reste à explorer afin d'améliorer la rigueur de la revue de littérature, notamment le modèle de mémoire utilisé pour le transport de la lumière dans \cite{seyb24from} et l'approche de convolution creuse plus récente de \textcite{xu25practical}. Ces ajouts fourniront une base théorique plus complète pour le travail d'implémentation à venir.
